\SourceLeaf{323}
\chapter*{ON THE THEORY OF BIQUADRATIC RESIDUES.}
\pdfbookmark[0]{On the Theory of Biquadratic Residues}{nachlass-biquadratic}
\renewcommand{\BookRunningTitle}{On the Theory of Biquadratic Residues}
\BookRule{.18\textwidth}
\begin{center}\textbf{[I.]}\end{center}
\pdfbookmark[1]{I}{nachlass-biquadratic-I}
\renewcommand{\Article}[1]{\par\medskip\Needspace{4\baselineskip}\pdfbookmark[2]{Art. #1}{nachlass-biquadratic-I-art-#1}\centerline{\textbf{#1.}}\nobreak\smallskip}
\Article{1}
We extend the domain of higher arithmetic by including complex quantities as well. In this investigation, we call every quantity $x+iy$ a complex integer when $x,y$ are real integers.

\Article{2}
The infinitely many complex integers are most conveniently represented visually by points in an unbounded plane; we call the point whose horizontal coordinate is $x$ and vertical coordinate $y$ the point $x+iy$; all points representing integers will be called integer points.

\Article{3}
To fix a definite convention, the horizontal coordinates shall always be positive to the left, and the vertical coordinates positive upwards.

\Article{4}
The straight line drawn from the point $x+iy$ to the point $x'+iy'$ will be called the straight line $(x+iy,x'+iy')$; wherever it matters, we also take its direction into account, and therefore distinguish the straight line $x+iy,x'+iy'$ from $x'+iy',x+iy$.

\SourceLeaf{324}
\Article{5}
For brevity, we shall also denote complex quantities such as $x+iy$ by a single letter, such as $z$.

\Article{6}
We call the figure bounded by the straight lines $zz',\allowbreak z'z'',\allowbreak z''z'''\ldots z^{n-1}z^n,\allowbreak z^nz$ the figure $zz'z''z'''\dotstwo z^n$. We do not exclude the case where some of these lines intersect.

\Article{7}
By $S(z,z',z''\ldots z^n)$ we generally denote a sum of as many real integers as there are integer points inside the figure, assigning the number $\pm1$, $\pm2$, $\pm3$, etc., to each point around which the boundary of the figure passes once, twice, three times, and so on; the upper signs apply when the boundary surrounds the point so that it lies on the right side of the figure, and the lower signs in the opposite case. Thus, if no sides of the figure intersect, $S(z,z',z''\dotstwo)$ is simply the number of points inside the figure, taken positively or negatively.

\Article{8}
It is evident that always
\[
\begin{aligned}
S(z,z',z''\ldots z^n)&=S(z',z'',z'''\dotstwo z^n,z)\\
&=S(z'',z'''\dotstwo z')\text{ etc.}\\
&=-S(z^n,z^{n-1}\dotstwo z'',z',z)\\
&=-S(z^{n-1}\pdot z^{n-2}\dotstwo z',z,z^n)\text{ etc.}
\end{aligned}
\]

\Article{9}
How points on the boundary itself are to be treated must still be specified more closely. In many cases no integer points can lie on the boundary at all; no convention is then needed. But if such points lie on the boundary $zz'$, inserting between $z$ and $z'$ a $+$ indicates that they are to be treated as if they lay to the right of the boundary, and a $-$ as if they lay to the left. We shall also sometimes insert a $0$ or $\frac12$ to indicate that they are to be counted either not at all or only at half their value on
\SourceLeaf{325}
each side. If any of the points $z,z',z''$, etc., is itself an integer point, it is not counted at all unless the contrary is expressly stated, except insofar as it may also occur as a point that is not a vertex.

\Article{10}
Theorems. If all $z,z',z''$, etc., are increased by the same integer, $S$ remains unchanged.

\Needspace{11\baselineskip}
If $i$ is replaced by $-i$ and every connecting sign by its opposite, $S$ changes only its sign.
\[
\begin{aligned}
S(z,z',z''\ldots z^n)={}&S(z,u,u'\dotstwo u^n,z',z'',z^n)\\
&-S(z,u,u'\dotstwo u^n,z')\\
={}&S(z,u,u',u''\dotstwo u^n,z^m,z^{m+1}\dotstwo z^n)\\
&-S(z,u,u',u''\dotstwo u^n,z^m,z^{m-1}\dotstwo z',z)
\end{aligned}
\]
Here the connecting signs must correspond, but must be reversed between the terms that run backwards.

If $\zeta$ is an integer $=a+bi$, then, when opposite connecting signs are opposite,
\[
S(z,z',z'+\zeta,z+\zeta)=[bx'-ay']-[bx-ay]
\]
Here, if $bx'-ay'$ is itself an integer, that integer must be taken for $[bx'-ay']$ when the connecting sign between $z'$ and $z'+\zeta$ is $+$, but $1$ or $\frac12$ less when this sign is $-$ or $\frac12$; the reverse applies to $bx-ay$.

The formula holds only when $a$ and $b$ have no common divisor; if their greatest common divisor is $=h$, the expression to use instead is
\[
h\left[\frac{bx'-ay'}h\right]-h\left[\frac{bx-ay}h\right]
\]
